% !TeX program = pdflatex
\documentclass[12pt,a4paper]{article}

% ==== Base packages ====
\usepackage[english]{babel}
\usepackage[utf8]{inputenc}
\usepackage[T1]{fontenc}
\usepackage[top=2.5cm,bottom=2.5cm,left=2.5cm,right=2.5cm]{geometry}
\usepackage{csquotes}
\usepackage{setspace}
\usepackage{microtype}
\usepackage{graphicx}
\usepackage{float}
\usepackage{booktabs}
\usepackage{array}
\usepackage{longtable}
\usepackage{tabularx}
\usepackage{enumitem}
\usepackage{titlesec}
\usepackage{hyperref}
\usepackage{xcolor}
\usepackage{caption}
\usepackage{listings}
\usepackage{fancyhdr}

% ==== Styles ====
\hypersetup{
  colorlinks=true,
  linkcolor=blue!60!black,
  urlcolor=blue!60!black,
  citecolor=blue!60!black,
  pdfauthor={MLR-X Team},
  pdftitle={End-User Operator Manual – MLR-X},
  pdfcreator={MLR-X}
}

\lstdefinestyle{mlrx}{
  basicstyle=\ttfamily\small,
  frame=single,
  breaklines=true,
  showstringspaces=false,
  columns=fullflexible
}

\setlist{nosep,leftmargin=2.1em}
\onehalfspacing
\pagestyle{fancy}
\fancyhf{}
\lhead{End-User Operator Manual — MLR-X 1.0}
\rhead{\thepage}

\titleformat{\section}{\Large\bfseries}{\thesection}{0.5em}{}
\titleformat{\subsection}{\large\bfseries}{\thesubsection}{0.5em}{}
\titleformat{\subsubsection}{\normalsize\bfseries}{\thesubsubsection}{0.5em}{}

% ==== Cover ====
\title{\vspace{-1.2cm}\textbf{MLR-X 1.0}\\End-User Operator Manual}
\author{MLR-X Team}
\date{\today}

\begin{document}
\maketitle

\begin{center}
\textit{This document describes the operational use of the MLR-X desktop application for data preparation, multiple linear regression (MLR) model search, validation, diagnostics, visualization, reporting, and command-line (CLI) automation.}
\end{center}

\tableofcontents
\newpage

% =========================================================
\section{Overview}
MLR-X is a desktop application (Tkinter) that guides users through preparing a \texttt{CSV} file, choosing the dependent variable, and filtering columns before launching an automatic search of multiple linear regression models. Everything is coordinated from the \textbf{Setup} tab, where you load data, set the delimiter, and \emph{start or cancel} analyses.

When a sweep finishes, the \textbf{Models} tab lists the combinations found (sorted by key metrics) and lets you \emph{reload saved results} to keep exploring without recomputation.

% --- Placeholder figure ---
\begin{figure}[H]
  \centering
  \fbox{\rule[0pt]{0pt}{120pt}\hspace{0.9\linewidth}}
  \caption{Main window overview: title \emph{``MLR-X 1.0''} (1000$\times$720 px) with tabbed notebook.}
  \label{fig:main-window}
\end{figure}

\subsection{Main tabs}
\begin{longtable}{@{}p{3cm}p{11.2cm}@{}}
\toprule
\textbf{Tab} & \textbf{Purpose} \\
\midrule
\endfirsthead
\toprule
\textbf{Tab} & \textbf{Purpose} \\
\midrule
\endhead
Setup &
Data loading and parameters: search method, sorting target, prediction clipping options, and saving configuration. Starting point of the workflow.\\
Models &
Results tables for training, internal and external validations (independent scrollable views). Metrics like R$^2$, RMSE, Q$^2$ per model.\\
Validation &
Choose which models to revalidate (top, multiple, all). Enable LOO or \emph{k}-fold. Run external validations on a different CSV; separate status lines per process.\\
Diagnostics &
Observation-level residuals and hat values, train/test filters, and downloadable correlation summaries to investigate outliers after selecting a specific model.\\
Visualization &
Preset plots (heatmaps, Predicted vs Actual, QQ, Williams, etc.), style controls, and legends to interpret the selected model visually.\\
Summary &
Generates a dossier-like report with global metrics, model information, and exportable equations.\\
Desktop (Variable Explorer) &
Reload data and results, choose variables for axes or model ranges, customize colors/markers/labels, and export figures.\\
\bottomrule
\end{longtable}

\subsection{Persistent state and performance}
\begin{itemize}
  \item MLR-X \textbf{retains} last results, analysis context, configurations, and derived metrics so specialized tabs can be reactivated without recalculation.
  \item An \textbf{event queue} and \textbf{worker thread} keep the UI responsive; progress, logs, and final results update as soon as they are ready.
  \item \textbf{Caches} for observations and correlations reduce waiting time when reopening diagnostics or visualizations within the same run.
\end{itemize}

% =========================================================
\section{System requirements and installation}
\subsection{Requirements}
\begin{itemize}
  \item Python environment with Tk/Tcl graphics support for Tkinter.
  \item Libraries: \texttt{NumPy}, \texttt{pandas}, \texttt{statsmodels}, \texttt{scikit-learn}, \texttt{Matplotlib}, \texttt{joblib}.
\end{itemize}

\subsection{Recommended installation}
\begin{enumerate}
  \item Create and activate a \textbf{virtual environment} (\texttt{python -m venv} or \texttt{conda}).
  \item Install dependencies:
\begin{lstlisting}[style=mlrx]
pip install numpy pandas statsmodels scikit-learn matplotlib joblib
\end{lstlisting}
  \item Verify Tkinter:
\begin{lstlisting}[style=mlrx]
python -m tkinter
\end{lstlisting}
\end{enumerate}

% =========================================================
\section{Quick start}
\subsection{Launch the GUI}
\begin{lstlisting}[style=mlrx]
python MLRX.py
\end{lstlisting}
The app creates an \texttt{MLRXApp} instance and enters the Tkinter main loop.

% =========================================================
\section{User interface tour}
\subsection{\textbf{Setup} tab}
\paragraph{\emph{Dataset} group}
\begin{enumerate}[label=\arabic*.]
  \item Specify the CSV (\emph{CSV file}) or click \emph{Browse}.
  \item \emph{Run analysis} starts the sweep; \emph{Preview data} shows a quick dataset summary; \emph{Cancel} stops execution.
  \item Set the \emph{CSV delimiter}, choose the \emph{dependent variable}, and mark \emph{non-predictor} columns.
  \item Enable \emph{Exclude near-constant predictors} and set the threshold (\%).
  \item Exclude rows using IDs/ranges in \emph{Exclude observations (IDs)} (e.g., \texttt{3,9,4-8}).
  \item Load/progress status appears in the upper-right corner of the box.
\end{enumerate}

\paragraph{\emph{Data splitting} group}
\begin{itemize}
  \item \emph{No split}: use the entire dataset.
  \item \emph{Random split}: define \emph{Test size (\%)}.
  \item \emph{Manual split}: lists/ranges of IDs for \textit{train}/\textit{test}.
\end{itemize}

\paragraph{\emph{Settings} group}
\begin{enumerate}[label=\arabic*.]
  \item Model search method: \textbf{EPR-S} or \textbf{All Subsets} (info via \emph{ℹ}).
  \item Key parameters: \emph{Max predictors}, \emph{Seeds}, \emph{Seed size}, statistical levels, \emph{export limits}, parallel jobs.
  \item Sorting target (R$^2$, adjusted R$^2$, or Q$^2$), \emph{covariance type}, prediction \emph{clipping}, and recommended buttons.
  \item Tooltips (ℹ) on every critical field.
\end{enumerate}

\subsection{Built-in validations and aids}
\begin{itemize}
  \item \emph{Preview data} reloads using the current configuration and previews train/test counts, response column, and first predictors; any read issue raises an error dialog.
  \item When starting analysis or saving a configuration, the app validates dataset paths, positive limits, integers, and warns on inverted clipping ranges.
  \item For splitting, the app enforces numeric 0–100 percentages, valid manual IDs, and an existing external file when chosen.
\end{itemize}

\subsection{Recommended practices}
\begin{itemize}
  \item \textbf{Seeds}: at least as many as predictors; use the \emph{Recommended} button.
  \item \textbf{Seed size}: around half of \emph{Max predictors}.
  \item \textbf{VIF threshold}: keep $<5$.
  \item \textbf{Correlation threshold}: 0–1; \textbf{Significance level}: 0.05 as a baseline.
\end{itemize}

% --- Placeholder figure ---
\begin{figure}[H]
  \centering
  \fbox{\rule[0pt]{0pt}{120pt}\hspace{0.9\linewidth}}
  \caption{\textbf{Setup} tab: \emph{Dataset}, \emph{Data splitting}, and \emph{Settings} groups.}
  \label{fig:setup}
\end{figure}

% =========================================================
\section{Running analyses and process management}
\subsection{Before you start}
\begin{enumerate}
  \item Load the dataset (\emph{Open dataset}) and verify the preview.
  \item Review delimiter, response, exclusions, validation choices, and filters.
  \item Choose the export destination (\texttt{models.csv} by default).
\end{enumerate}

\subsection{Start the analysis}
\begin{enumerate}
  \item Ensure no other run is active.
  \item Click \emph{Run analysis} and confirm the results file path.
  \item The app validates the destination and reloads the dataset using your settings.
  \item Status switches to \emph{Running}, prior results are cleared, and the log starts.
\end{enumerate}

\subsection{Progress, cancellation, incidents}
\begin{itemize}
  \item \textbf{Progress}: bottom-right bar (\%) and \emph{Log} with evaluated models and events.
  \item \textbf{Cancel}: switches to \emph{Cancelling…} while jobs are gracefully stopped.
  \item \textbf{Errors}: dialog with cause and overall \emph{Error} state.
\end{itemize}

\subsection{Results and metadata}
\begin{enumerate}
  \item The CSV is written to the confirmed path (\texttt{models.csv} by default).
  \item First line: a JSON metadata block between \verb|#EPRS-S ... ;#|.
  \item The last confirmed export path is kept as the default for future runs.
\end{enumerate}

% =========================================================
\section{Understanding the \textbf{Models} tab}
The \textbf{Models} (or \textbf{Results}) tab groups three main tables and an \textit{observation diagnostics} view. They share the model ID and predictor list; each focuses on different metrics.

\subsection{Top models (training)}
\textbf{Columns}: \emph{Model}, \emph{Predictors}, \emph{N pred}, R$^2$ (train), RMSE (train), $s$ (train), MAE (train), adjusted R$^2$, VIF max, VIF avg.\\
\textbf{How to read}: balance R$^2$/adj-R$^2$; compare RMSE/MAE/$s$; check VIF; watch \emph{N pred}.\\
\textbf{Relation}: selection feeds \textbf{Validation} and \textbf{Diagnostics}.

\subsection{Internal validation (cross-validation)}
\textbf{Columns}: \emph{Model}, \emph{Predictors}, \emph{N pred}, Q$^2$/RMSE/MAE/$s$ for LOO and for \emph{k}-fold.\\
\textbf{Relation}: launched from \textbf{Validation}; results populate this table.

\subsection{External validation}
\textbf{Columns}: \emph{Model}, \emph{Predictors}, \emph{N pred}, Q$^2_\mathrm{F3}$, RMSE (ext), MAE (ext), Q$^2_\mathrm{F2}$, Q$^2_\mathrm{F1}$.

\subsection{Observation diagnostics}
\textbf{Columns}: \emph{Observation}, \emph{Set}, \emph{Actual}, \emph{Predicted}, \emph{Predicted (LOO)}, \emph{Residual}, \emph{Residual (LOO)}, \emph{Z value}, \emph{Hat}, \emph{Std. residual}, \emph{Std. residual by LOO}.

\subsection{Sorting and refresh}
\begin{itemize}
  \item \textbf{Sorted by} syncs ordering across all tables.
  \item \textbf{Max rows} caps visible/exported rows (default up to 5000).
  \item Any action affecting metrics (sorting, added validations, external evaluation) \emph{synchronizes} \textbf{Results}, \textbf{Summary}, \textbf{Validation}, and \textbf{Diagnostics}.
\end{itemize}

% --- Placeholder figure ---
\begin{figure}[H]
  \centering
  \fbox{\rule[0pt]{0pt}{140pt}\hspace{0.9\linewidth}}
  \caption{\textbf{Models} tab: training, internal, and external validation tables.}
  \label{fig:models}
\end{figure}

% =========================================================
\section{Internal and external validation (\textbf{Validation})}
\subsection{Model selection and LOO/\emph{k}-fold}
\begin{itemize}
  \item \textbf{Model selection}: \emph{Top models} (with limit), \emph{Multiple models} (lists/ranges: \texttt{1-5,7,9}), or \emph{All models}.
  \item \textbf{Internal validation}: LOO (default) and \emph{k}-fold (enables \emph{Folds} and \emph{Repeats}).
  \item Status displayed next to \emph{Run internal validation}.
\end{itemize}

\subsection{Execution and interpretation}
\begin{itemize}
  \item Requires models in memory.
  \item Validates that at least one method is selected and integer fields are plausible (\emph{folds} $\ge 2$, \emph{repeats} $\ge 1$, \emph{folds} $\le n$).
  \item The log reports how many models and which schemes; on finish, it summarizes totals and warnings.
\end{itemize}

\subsection{External validation flow}
\begin{enumerate}
  \item Provide the holdout CSV and its delimiter (\verb|;|, \verb|,|, tab, \verb/|/).
  \item Checks: non-empty file, same response column as training.
  \item Choose models as in internal validation; the \emph{External validation} status summarizes results.
\end{enumerate}

% =========================================================
\section{\textbf{Summary} tab and reporting}
\subsection{Layout}
Top \emph{Model selection} block (model list, status, \emph{Use recommended covariance}, \emph{Export Summary…}) and a main pane with a full-width, scrollable text box.

\subsection{Content and metrics}
\begin{itemize}
  \item Re-estimates OLS on model change, reports \emph{covariance type}, and rebuilds the coefficients table (SE, $t$, $p$, CIs).
  \item \textbf{Training}: R$^2$, adj-R$^2$, RMSE, MAE, $s$, VIF max, predictor correlations.
  \item \textbf{Internal}: LOO/\emph{k}-fold averages (with repetition notes).
  \item \textbf{External}: Q$^2_\mathrm{F1,F2,F3}$, RMSE, MAE.
  \item \textbf{Advanced}: $F$ and $p$, log-likelihood, AIC, BIC, Omnibus/JB, skewness, kurtosis, Durbin–Watson, condition number, Lin’s concordance coefficient, etc.
\end{itemize}

\subsection{Recommended covariance}
\begin{itemize}
  \item \emph{Use recommended covariance} runs BP/White and leverage checks; if heteroskedasticity or high leverage is detected, applies a robust estimator (HC3 by default).
  \item \emph{Covariance Diagnostics} documents tests, suggested/applied type, and rationale.
\end{itemize}

\subsection{Export}
\emph{Export Summary…} saves the entire text content to \texttt{.txt} (file name includes the model ID).

% --- Placeholder figure ---
\begin{figure}[H]
  \centering
  \fbox{\rule[0pt]{0pt}{120pt}\hspace{0.9\linewidth}}
  \caption{\textbf{Summary} tab: example dossier with metrics and equation.}
  \label{fig:summary}
\end{figure}

% =========================================================
\section{Diagnostics for observations and correlations}
\subsection{Selection and dataset filters}
\begin{itemize}
  \item \textbf{Setup}: choose dataset and exclusions (IDs, near-constant) before running.
  \item \textbf{Data splitting}: \emph{none}/\emph{random}/\emph{manual}/\emph{external}.
  \item \textbf{Diagnostics}: \emph{Dataset} filter (Training/Testing/Both) with holdout validation.
  \item \textbf{Visualization}: shares dataset selector and validation logic.
  \item Automatic propagation of the \emph{holdout ready} state across tabs.
\end{itemize}

\subsection{Diagnostic columns}
\begin{itemize}
  \item \textbf{Observation/Set}: normalized ID and membership.
  \item \textbf{Actual vs Predicted} and \textbf{Predicted (LOO)}.
  \item \textbf{Residual} and \textbf{Residual (LOO)}.
  \item \textbf{Z value} (studentized), \textbf{Hat} (guide threshold $3(p+1)/n$).
  \item \textbf{Std. residual} / \textbf{Std. residual by LOO}.
  \item \textbf{Cook’s distance} / \textbf{Cook’s distance LOO}.
\end{itemize}

\subsection{Exports and correlations}
\begin{itemize}
  \item \emph{Export CSV} enabled only when data are visible; respects the dataset filter.
  \item Correlation matrix: adds a \emph{Variable} column and uses the lower triangle (no duplicates).
  \item \emph{Include dependent variable}: enabled only when the extended matrix is ready.
\end{itemize}

% =========================================================
\section{Variable exploration and advanced visualizations}
\subsection{Variable Explorer (Desktop)}
\paragraph{Loading}
\begin{itemize}
  \item \textbf{CSV file / Browse / Load}: pick and load a data file; the file name and load status appear below.
  \item \textbf{Results file / Browse / Load}: load prior model results to unlock extra features.
\end{itemize}

\paragraph{Axes and scope}
\begin{itemize}
  \item \emph{All variables} or \emph{Only variables of model} (type model ID).
  \item \emph{X axis} / \emph{Y axis} lists update based on the chosen scope.
\end{itemize}

\paragraph{Legend, labels, references}
\begin{itemize}
  \item \emph{Show legend} and \emph{Legend location}.
  \item \emph{Label mode} (ID, name, etc.); left-click to toggle labels and drag them; right-click clears all labels.
  \item \emph{Show linear fit} (optional R$^2$), \emph{Show gridline}, \emph{Partial residuals}.
\end{itemize}

\paragraph{Appearance and export}
\begin{itemize}
  \item \emph{Marker size/color/style}, \emph{X/Y min–max}, \emph{Tick step}, \emph{Aspect ratio} (1, 4:3, 16:9, 9:16).
  \item \emph{Set axis labels…} and \emph{Save current plot} (PNG/TIFF/PDF/SVG).
\end{itemize}

\subsection{Visualization — plot types}
\begin{longtable}{@{}p{6.2cm}p{8.2cm}@{}}
\toprule
\textbf{Name in app} & \textbf{What it shows} \\
\midrule
Correlation heatmap & Correlations among model predictors (redundancy). \\
Correlation heatmap (with dependent) & Adds the target variable to the heatmap. \\
Predicted vs Actual / by LOO & Predictions vs observations (standard and LOO). \\
Residual vs Actual / by LOO & Residuals vs observed values. \\
Residuals vs Predictions / by LOO & Residuals vs predictions to spot patterns. \\
Scale-Location / by LOO & Homoskedasticity via standardized residuals. \\
Q-Q residuals / by LOO & QQ-plot vs normal (standard/LOO). \\
Residual distribution / by LOO & Histogram/density of standardized residuals. \\
Williams plot / by LOO & Studentized residuals vs leverage with $h^*$ cutoffs. \\
Cook's distance vs leverage / by LOO & Influential observations (Cook + leverage). \\
Hat-values plot & Leverage per observation and thresholds. \\
Predicted vs Hat & Leverage vs predictions. \\
\bottomrule
\end{longtable}

\paragraph{Common controls}
\begin{itemize}
  \item \emph{Show legend} / \emph{Legend location}, \emph{Label mode}, \emph{Show identity line} (e.g., in Predicted vs Actual), $h^*$ line color, \emph{Show linear fit}, \emph{Show gridline}, separate styles for training/testing, synchronized ranges, and export.
\end{itemize}

% --- Placeholder figure ---
\begin{figure}[H]
  \centering
  \fbox{\rule[0pt]{0pt}{140pt}\hspace{0.9\linewidth}}
  \caption{Visualization examples: Predicted vs Actual, QQ residuals, Williams plot.}
  \label{fig:viz}
\end{figure}

% =========================================================
\section{Model search algorithms}
\subsection{EPR-S (Expand–Perturb–Reduce–Swap)}
\begin{enumerate}
  \item \textbf{Objective and start}: choose the target metric (e.g., R$^2$) and an initial predictor set with a size cap.
  \item \textbf{Expand}: try adding non-included predictors that improve the metric without exceeding \emph{Max predictors}.
  \item \textbf{Perturb}: 1--1 swaps to escape local optima; keep only improvements.
  \item \textbf{Reduce}: remove predictors when equal or better quality is achieved (parsimony).
  \item \textbf{Correlation cleanup}: if $|r|$ exceeds the threshold, drop the less useful variable.
  \item \textbf{Swap}: controlled re-introductions when they bring improvement; converge when no further gains.
  \item \textbf{Quality rules}: VIF and \emph{Minimum report R$^2$} filters; unique models kept and sorted.
\end{enumerate}

\subsection{VIF, reporting threshold, progress}
\begin{itemize}
  \item \textbf{VIF threshold}: set in \textbf{Setup} $\rightarrow$ \emph{Settings} (recommended $<5$).
  \item \textbf{Minimum report R$^2$ (or Q$^2$)}: models below the threshold are not added.
  \item \textbf{Progress}: seed counter and percentage bar; each completed seed advances the bar.
\end{itemize}

\subsection{All Subsets}
\begin{itemize}
  \item Generates/evaluates all combinations up to \emph{Max predictors}; exact but combinatorial.
  \item Recommended for $\lesssim25$ predictors; shares VIF/threshold filters with EPR-S.
\end{itemize}

\subsection{Recommendations}
\begin{itemize}
  \item Use \textbf{All Subsets} for small predictor sets when exhaustiveness matters.
  \item Use \textbf{EPR-S} for larger spaces or when you need speed with quality controls.
\end{itemize}

% =========================================================
\section{Results format, metadata, and exports}
\subsection{Main results CSV}
\begin{itemize}
  \item \texttt{models.csv} (separator \verb|;|).
  \item First line: JSON metadata between \verb|#EPRS-S| and \verb|;#|.
  \item Rows: model ID, variables, and metrics (R$^2$, RMSE, MAE, $s$, LOO, \emph{k}-fold, external when available) with four decimals.
\end{itemize}

\subsection{Paths, permissions, reuse}
\begin{itemize}
  \item Default: app folder; you can choose a custom destination.
  \item App checks permissions and file locks.
  \item \textbf{Variable Explorer} $\rightarrow$ \emph{Results file / Load} restores a session from the JSON header.
\end{itemize}

\subsection{Complementary exports}
\paragraph{Filtered datasets}
\begin{itemize}
  \item \emph{Create filtered dataset} wizard: choose model or range; preview of included columns (IDs, auxiliaries, target, predictors).
  \item Suggested name \texttt{model\_\textless id\textgreater\_dataset.csv}; uses the \verb|;| separator.
\end{itemize}

\paragraph{Figures}
\begin{itemize}
  \item \emph{Save current plot}: PNG/TIFF/PDF/SVG with current on-screen configuration.
  \item Remembers last format; warns if no data are loaded.
\end{itemize}

% =========================================================
\section{CLI usage and automation}
\subsection{\texttt{.conf} file format}
UTF-8 text with \emph{Label = value} lines, blanks, and \verb|#| comments. Main blocks:
\begin{itemize}
  \item \textbf{Dataset path}, \textbf{Delimiter options} (choose among \verb|;|, \verb|,|, tab, \verb/|/).
  \item \textbf{Dependent variable} and \textbf{Non-variable columns}.
  \item \textbf{Exclude constant columns}, \textbf{Constant threshold}, \textbf{Excluded observation IDs}.
  \item \textbf{Search method}, \textbf{Covariance type}, \textbf{Target metric}.
  \item \textbf{Maximum variables}, \textbf{Top models to export}, \textbf{Parallel jobs}.
  \item \textbf{Split mode}: \emph{none}/\emph{random}/\emph{manual}/\emph{external} (with associated fields).
  \item \textbf{Output path}.
\end{itemize}

\subsection{Execution}
\begin{lstlisting}[style=mlrx]
python MLRX.py path/to/config.conf
\end{lstlisting}

During execution:
\begin{enumerate}
  \item Config is confirmed and validated; applied to the dataset.
  \item Console shows progress (\%) and phases.
  \item On finish: total models, total time, and whether valid results were found.
  \item If models exist, the CSV is exported to \textbf{Output path}.
\end{enumerate}

\subsection{Automation}
\begin{itemize}
  \item \textbf{Scheduling}: use \texttt{cron}, Windows Task Scheduler, or shell scripts.
  \item \textbf{Consistent outputs}: the CSV overwrites; append date/time in a wrapper script for versioning.
  \item \textbf{Pipeline integration}: the first CSV line summarizes config for easy downstream parsing.
  \item \textbf{Monitoring}: redirect \emph{stdout/stderr} to log files.
\end{itemize}

% =========================================================
\section{EPR-MLR quick guide (CSV prep and splitting)}
\subsection{Prepare the CSV}
\paragraph{Minimum content}
\begin{itemize}
  \item An ID column (e.g., \emph{ID} or \emph{Sample}).
  \item At least one numeric predictor.
  \item The target (dependent) variable.
\end{itemize}
\paragraph{Format}
\begin{itemize}
  \item CSV (UTF-8) with a consistent separator (comma, semicolon, tab).
\end{itemize}
\paragraph{Special columns}
\begin{itemize}
  \item Mark \emph{non-variables} (IDs, dates, etc.).
  \item If absent, the app creates a sequential numeric ID on import.
\end{itemize}
\paragraph{Missing values and text}
\begin{itemize}
  \item Address missing values; convert categorical text to numeric codes if used as predictors.
\end{itemize}
\paragraph{Tip} Validate headers and types in a spreadsheet before importing.

\subsection{Configure import in the app}
\begin{enumerate}
  \item Open the app and click \emph{Load data}.
  \item Choose the CSV; set separator and encoding if needed.
  \item Select the dependent variable.
  \item Provide indices/names of non-predictor columns.
  \item Use \emph{Preview} to confirm interpretation.
\end{enumerate}

\subsection{Preprocessing options}
\begin{itemize}
  \item \emph{Exclude observations}: IDs and ranges.
  \item \emph{Remove (near-)constants}: tolerance threshold (\% of most frequent value).
  \item \emph{Convert non-numeric}: columns that cannot be converted are dropped from predictors.
\end{itemize}

\subsection{Choose a splitting mode}
\begin{center}
\begin{tabularx}{\linewidth}{@{}>{\bfseries}l X X@{}}
\toprule
Mode & When to use & Requirements \\
\midrule
None & Train on full dataset. & No extra configuration.\\
Random & Random train/test split. & Define \emph{test size} (e.g., 0.2).\\
Manual & Full control of train/test. & Lists/ranges of non-overlapping IDs.\\
External & External holdout file. & Path, delimiter, same columns as training CSV.\\
\bottomrule
\end{tabularx}
\end{center}

\subsection{Good practices before training}
\begin{itemize}
  \item Document exclusions and keep a clean copy of the original CSV.
  \item Review columns dropped for non-numeric content.
  \item For external holdout, ensure identical headers and order.
\end{itemize}

% =========================================================
\section{Saving configurations and reuse}
\subsection{Save \texttt{.conf} from the GUI}
Use the GUI to save the current configuration (\texttt{write\_configuration\_file}) after confirming a path in the save dialog.

\subsection{\texttt{.conf} content}
Stores data paths, delimiter, dependent variable choice, constant filters, search method, target metric, export limits, clipping options, and splitting settings (random, manual, external). In CLI, \texttt{parse\_configuration\_file} validates required fields and indices.

% --- Placeholder figure ---
\begin{figure}[H]
  \centering
  \fbox{\rule[0pt]{0pt}{110pt}\hspace{0.9\linewidth}}
  \caption{Saving a configuration and sample \texttt{.conf} layout.}
  \label{fig:conf}
\end{figure}

% =========================================================
\section{Operational best practices}
\begin{itemize}
  \item Review parameters before each run and align with dataset structure and project goals.
  \item Save configurations to enable batch/CLI reproducibility.
  \item Use \textbf{Diagnostics} and \textbf{Validation} to confirm robustness before decisions.
\end{itemize}

% =========================================================
\section{Quick troubleshooting (operational FAQ)}
\begin{itemize}
  \item \textbf{GUI does not open}: verify \texttt{python -m tkinter}. Ensure Tk/Tcl is available in your environment.
  \item \textbf{CSV fails to load}: check path, permissions, and delimiter; try \emph{Preview data}.
  \item \textbf{No models exported}: review \emph{Minimum report R$^2$}/Q$^2$ and \emph{VIF threshold}.
  \item \textbf{External validation fails}: the target column must match; CSV must include the same predictors.
  \item \textbf{Empty plot}: select a valid model/dataset or load results (\emph{Results file}).
\end{itemize}

% =========================================================
\section{Glossary}
\begin{description}[style=unboxed,leftmargin=0cm,labelsep=0.5em,font=\normalfont\bfseries]
  \item[R$^2$/adj-R$^2$] Coefficient of determination (adjusted accounts for number of predictors).
  \item[Q$^2$] Out-of-sample predictive power (LOO/\emph{k}-fold).
  \item[RMSE/MAE/$s$] Root mean squared error, mean absolute error, and model standard error.
  \item[VIF] Variance Inflation Factor; multicollinearity (recommended $<5$).
  \item[Hat] Leverage; rule-of-thumb threshold $3(p+1)/n$.
  \item[Williams plot] Studentized residuals vs leverage with critical $h^*$.
  \item[HC0–HC3] Heteroskedasticity-robust covariance estimators.
\end{description}

% =========================================================
\section{Command reference (summary)}
\subsection*{GUI}
\begin{lstlisting}[style=mlrx]
python MLRX.py
\end{lstlisting}

\subsection*{CLI}
\begin{lstlisting}[style=mlrx]
python MLRX.py path/to/config.conf
\end{lstlisting}

% --- Final placeholder figure ---
\begin{figure}[H]
  \centering
  \fbox{\rule[0pt]{0pt}{110pt}\hspace{0.9\linewidth}}
  \caption{Example \textbf{Validation} tab with internal/external status indicators.}
  \label{fig:validation}
\end{figure}

\end{document}
